\section{Separate powers of the target poles}
\label{sec:split}

Choose an analytic square root \(s\) of \(D\) on a local chart where
\(D\ne0\). Write \(z_\pm=c\pm s\), \(d_0=2s\), \(x_j=w_j-s\), and
\(y_j=w_j+s=x_j+d_0\). The split-pole moments are
\begin{equation}
 P_{p,q}=\sum_j x_j^{-p}y_j^{-q},
 \qquad p,q\in\mathbb Z_{\geq0},\quad p+q\geq2.
 \label{split:def}
\end{equation}
Both \(\lvert x_j\rvert\) and \(\lvert y_j\rvert\) are at least
\(3\lvert w_j\rvert/4\) by \eqref{eq:domain}. Thus all of
\eqref{split:def}, including its pure powers, converge absolutely.

The two moment families have an exact algebraic relation.
For \(h=\max(p,q)\) one has
\begin{equation}
 \frac1{(w-s)^p(w+s)^q}
    =\frac{(w-s)^{h-p}(w+s)^{h-q}}{(w^2-D)^h}.
 \label{split:reduction}
\end{equation}
The numerator has degree at most \(2h-2\) in \(w\).
Write an even power as \(w^{2k}=(Q+D)^k\) and an odd power as
\(w^{2k+1}=w(Q+D)^k\), with \(Q=w^2-D\).
Then \(k\leq h-1\) in the even terms and \(k\leq h-2\) in the
odd terms. Expansion of the binomial gives only \(S_n\) with
\(n\geq1\) and \(T_n\) with \(n\geq2\), with polynomial
coefficients in \(s\) after \(D=s^2\). Conversely,
\begin{equation}
 S_n=P_{n,n}\quad(n\geq1),\qquad
 T_n=\tfrac12(P_{n-1,n}+P_{n,n-1})\quad(n\geq2).
 \label{split:inverse}
\end{equation}
In particular, since $w=(w-s)+s=(w+s)-s$,
\begin{equation}
 T_2=P_{1,2}+sP_{2,2}=P_{2,1}-sP_{2,2}.
 \label{split:T2}
\end{equation}
This uses \(xy=Q\) and \(x+y=2w\). Both changes of expressions
use the absolutely convergent index ranges.

\begin{proposition}[Split-pole recurrence]\label{split:prop}
For \(m,n\geq1\) write
\begin{equation}
 C_{m,n}=\sum_{i=2}^{m}P_{i,0}P_{m+1-i,n}
            +\sum_{j=1}^{n-1}P_{m,j}P_{0,n+1-j}.
 \label{split:convolution}
\end{equation}
On the axes use \(C_{m,0}=\sum_{i=2}^{m-1}P_{i,0}P_{m+1-i,0}\)
and \(C_{0,n}=\sum_{j=2}^{n-1}P_{0,j}P_{0,n+1-j}\) for the
respective index at least two. Let \(A_d=2/d_0-d_0M\).
For \(p,q\geq1\), the complete hierarchy \cite{BakerCM08} is
\begin{align}
 P_{p,q}'={}&p\bigl(C_{p+1,q}-(p+3)P_{p+2,q}-(q+2)P_{p+1,q+1}
                                  +A_dP_{p+1,q}\bigr)\notag\\
   &+q\bigl(C_{p,q+1}-(p+2)P_{p+1,q+1}-(q+3)P_{p,q+2}
                                  -A_dP_{p,q+1}\bigr).
 \label{split:mixed}
\end{align}
The pure powers satisfy
\begin{align}
 P_{p,0}'&=p\bigl(C_{p+1,0}-(p+3)P_{p+2,0}-2P_{p+1,1}
                                   +(2/d_0)P_{p+1,0}\bigr),
       \qquad p\geq2,\notag\\
 P_{0,q}'&=q\bigl(C_{0,q+1}-2P_{1,q+1}-(q+3)P_{0,q+2}
                                   -(2/d_0)P_{0,q+1}\bigr),
       \qquad q\geq2 .
 \label{split:axes}
\end{align}
The separation obeys \(d_0'=4/d_0-2d_0M\). The equations
are valid on the separated charts, including exterior multiplicities.
\end{proposition}

\begin{proof}
At a finite simple-root slice define \(P_{1,0}\) and \(P_{0,1}\)
as finite sums. The target velocities are
\[
 z_+'=2/d_0-2P_{1,0},\qquad z_-'=-2/d_0-2P_{0,1}.
\]
The identity \(x_j^{-1}-y_j^{-1}=d_0/Q_j\) gives
\(P_{1,0}-P_{0,1}=d_0M\) and the stated equation for \(d_0'\).
For the complete divisor this difference is the absolutely
convergent series \(\sum_j d_0/Q_j\), or the limit of differences
with a common cutoff.

At an exterior root the contribution of the two targets to
\(\rho_j'\) is \(2(x_j^{-1}+y_j^{-1})\).
Subtracting \(z_+'\) for \(x_j'\) and \(z_-'\) for \(y_j'\), and
differentiating the weight \(f_{p,q}(w)=(w-s)^{-p}(w+s)^{-q}\),
gives the following two contributions to \(P_{p,q}'\):
\begin{align*}
 &p\bigl((2/d_0-2P_{1,0})P_{p+1,q}
                    -2P_{p+2,q}-2P_{p+1,q+1}+U_{p+1,q}\bigr),\\
 &q\bigl((-2/d_0-2P_{0,1})P_{p,q+1}
                    -2P_{p+1,q+1}-2P_{p,q+2}+U_{p,q+1}\bigr).
\end{align*}
Omit a block with zero prefactor. Here the exterior pair term is
\[
 U_{m,n}=-2\sum_{j<k}
        \frac{f_{m,n}(w_j)-f_{m,n}(w_k)}{w_j-w_k}.
\]
Splitting the difference between the two factors gives
\begin{equation}
 -\frac{f_{m,n}(u)-f_{m,n}(v)}{u-v}
   =\sum_{i=1}^{m} f_{i,0}(u)f_{m+1-i,n}(v)
     +\sum_{j=1}^{n} f_{m,j}(u)f_{0,n+1-j}(v).
 \label{split:divdiff}
\end{equation}
For the first sum subtract \((u-s)^{-m}\) and \((v-s)^{-m}\)
using the difference of reciprocal powers, with
\((v+s)^{-n}\) held fixed. For the second use the corresponding
reciprocal powers at \(-s\) with \((u-s)^{-m}\) fixed.
The full kernel is symmetric in \(u,v\). Therefore \(U_{m,n}\)
is the sum of the right side over ordered off-diagonal pairs.
Subtracting the diagonal gives \(-mP_{m+1,n}\) and
\(-nP_{m,n+1}\) in the two sums.

When \(m,n\geq1\) the term \(i=1\) and the term \(j=n\) yield
\((P_{1,0}+P_{0,1})P_{m,n}\). Their combination with the two
moving-target blocks is \(-d_0MP_{p+1,q}\) and
\(+d_0MP_{p,q+1}\), respectively. The other products form
\eqref{split:convolution}. Adding the two target terms in each
block to the diagonal subtractions gives the four coefficients
in \eqref{split:mixed}. On the \(m\)-axis the two endpoints
\(i=1,m\) give \(2P_{1,0}P_{m,0}\), and on the \(n\)-axis the
endpoints \(j=1,n\) give \(2P_{0,1}P_{0,n}\). They cancel the
first-power sums in the corresponding blocks. The diagonals
here are \(-mP_{m+1,0}\) and \(-nP_{0,n+1}\). This proves
\eqref{split:axes}. The cancellation precedes all limiting steps.

The fixed physical circle \(\Gamma\) of Proposition~\ref{prop:trace}
also gives the contour formula for the split weight
\(k(z)=(z-z_+)^{-p}(z-z_-)^{-q}\).
The integral of \(k\) vanishes by its decay \(O(z^{-p-q})\).
The residue of \(k(z)/(z-\rho_j)\) at an exterior pole is
\(k(\rho_j)\), so the integral on \(\Gamma\) gives its negative.
In the convergent product logarithmic derivative each constant
regularizer multiplies the zero integral of \(k\). Hence
\begin{equation}
 P_{p,q}=-\frac1{2\pi\ii}\oint_\Gamma
                  k(z)\left(\frac{H_z}{H}-\frac{Q_{*,z}}{Q_*}\right)
                   \,\dd z.
 \label{split:trace}
\end{equation}
The \(Q_{*,z}/Q_*\) term accounts for both target poles.
Its weighted integral vanishes since all poles of \(kQ_{*,z}/Q_*\)
are enclosed and the decay is \(O(z^{-p-q-1})\).
For \(p=2,q=0\) the individual target residues are
\(-d_0^{-2}\) at \(z_+\) and \(+d_0^{-2}\) at \(z_-\).
Separating the two pole contours would retain these two terms.
For \(p=q\) the target term is the spatial derivative of
\(-Q_*^{-p}/p\).

Differentiate \eqref{split:trace} on the fixed circle, including
\(k_t=k(p z_+'/(z-z_+)+q z_-'/ (z-z_-))\).
The first contour jets converge by the argument of
Lemma~\ref{lem:slice-passage}. At a finite multiple-root slice
the root factors can be perturbed without changing the separated
counts, then evolved from that slice. The derivatives of \(z_\pm\)
come from the jets of \(c,D\) and \(s'=D'/(2s)\), with \(s\) bounded
away from zero on each compact subchart. This also applies to
the opposite-paired canonical product truncations there.
The omitted occurrences in \eqref{split:def} contribute
\(O(N^{1-p-q}\log N)\) at fixed orders, by the strip and unit count
as in the proof of Theorem~\ref{thm:hierarchy}. Every term of the
already canceled equations converges. This proves the complete
laws. At an interval endpoint contained in the chart the derivatives
are inward.
\end{proof}

For example, on the diagonal the identities
\[
 P_{n+1,n}-P_{n,n+1}=d_0 S_{n+1},\qquad
 P_{n+2,n}+P_{n,n+2}=2S_{n+1}+d_0^2S_{n+2}
\]
give, directly from \eqref{split:mixed},
\begin{align*}
 S_n'={}&n(C_{n+1,n}+C_{n,n+1})-4n(n+2)S_{n+1}\\
        &\quad-n(n+3)d_0^2S_{n+2}-nd_0^2MS_{n+1}.
\end{align*}
At \(n=1\) the convolution is \(M(2M+d_0^2 S_2)\).
At \(n=2\) its five products are
\[
 S_2(P_{2,0}+P_{0,2})+P_{3,0}P_{1,2}+P_{2,1}P_{0,3}
                       +P_{3,1}P_{0,2}+P_{2,0}P_{1,3}.
\]
Pair these products using \(x+y=2w\), \(y-x=d_0\), and \(xy=Q\).
For the two crossed pairs use
\begin{align*}
 P_{3,0}+P_{0,3}&=2T_2+2d_0^2T_3,&
 P_{3,0}-P_{0,3}&=d_0(3S_2+d_0^2S_3),\\
 P_{3,1}+P_{1,3}&=2S_2+d_0^2S_3,&
 P_{3,1}-P_{1,3}&=2d_0 T_3,\\
 P_{2,0}+P_{0,2}&=2M+d_0^2 S_2,&
 P_{2,0}-P_{0,2}&=2d_0 T_2,\\
 P_{2,1}+P_{1,2}&=2T_2,&
 P_{2,1}-P_{1,2}&=d_0 S_2 .
\end{align*}
The terms \(T_2T_3\) and \(d_0^4S_2S_3\) cancel, leaving
\(4MS_2+2T_2^2+d_0^2MS_3+(d_0^2/2)S_2^2\).
The two values of \(n\) reproduce \eqref{eq:low-orders} for \(S_1,S_2\).

\subsection{Coalescing labels and an exact quartic}

Changing \(s\) to \(-s\) swaps the two indices, and the right
sides of the recurrence transform by the same exchange. At real
times, conjugation fixes the moments for \(D>0\) and swaps the
two indices for \(D<0\). At a merger the two raw labels have the
same limit. In fact, \eqref{split:reduction} and the analytic
rational moments give a decomposition
\[
 P_{p,q}(t)=A(t)+s(t)B(t)
\]
with \(A,B\) locally analytic in \(t\). They are obtained by
separating the even and odd polynomial powers of \(s\) and
substituting \(s^2=D(t)\). At a target merger \(t_*\), expansion of
the summand in \(s\) at the nonzero exterior displacements gives
\[
 A(t_*)=\sum_j w_j(t_*)^{-p-q},\qquad
 B(t_*)=(p-q)\sum_j w_j(t_*)^{-p-q-1}.
\]
Both sums converge absolutely. The symmetric combination and
the divided antisymmetric combination recover \(A\) and \(B\).
An individual coordinate has derivative asymptotic to \(B(t_*)/s\)
if \(B(t_*)\ne0\), since \eqref{eq:pairlaw} gives \(D'(t_*)=2\).
The condition on \(B(t_*)\) remains open at a canonical merger.

The singularity occurs in a finite even quartic heat field.
For \(b>0\) consider
\begin{equation}
 F(t,z)=(z^2-b^2)^2-t(12z^2-4b^2)+12t^2 .
 \label{split:quartic}
\end{equation}
It satisfies \(F_t=-F_{zz}\) by differentiation. A fixed small
circle around \(+b\) counts the two targets near \(t=0\). The
exterior pair is their reflection, with center \(-c\) and the same
\(D\). Setting \(u=z-b\) at \(t=0\), one has
\[
 \frac{F_t}{F}
  =\frac{-8b^2-24bu-12u^2}{u^2(2b+u)^2}
  =-2u^{-2}-\frac4b u^{-1}+O(1).
\]
Comparison with the derivative of the target quadratic gives
\(D'(0)=2\), \(c'(0)=2/b\). The mean exterior displacement
\(\bar w=-2c\) thus has value \(u_0=-2b\) and derivative \(8/u_0\).
For a weight \(f(w,D)\) analytic near \((u_0,0)\), the trace over
the exterior pair is
\[
 f(\bar w+s,D)+f(\bar w-s,D)
    =2f(\bar w,D)+Df_{ww}(\bar w,D)+O(D^2).
\]
The derivative at the double is
\((16/u_0) f_w+4f_D+2f_{ww}\), evaluated at \((u_0,0)\).
For \(f=Q^{-n}\), the three contributions have coefficients
\(-32n,4n,4n(2n+1)\) of \(u_0^{-2n-2}\). For \(f=wQ^{-n}\)
they have coefficients \(16(1-2n),4n,-4n(1-2n)\) of
\(u_0^{-2n-1}\). In particular
\[
 S_n'(0)=8n(n-3)u_0^{-2n-2},\qquad
 T_n'(0)=8(n^2-4n+2)u_0^{-2n-1}.
\]
Independently, inserting \(S_n(0)=2u_0^{-2n}\),
\(T_n(0)=2u_0^{1-2n}\), and \(D(0)=0\) in the regular equations
gives the coefficients
\[
 8n^2+8n(n-1)-8n(n+2),
 \qquad
 8(2n-1)(n-1)-8(n^2+n-1),
\]
respectively. These reduce to the same derivatives.

Here \(P_{2,0}=S_1+2DS_2+2sT_2\) by \eqref{split:reduction},
and \(T_2(0)=2u_0^{-3}\ne0\). Thus the square-root term is
nonzero on either label for \(t>0\) near zero, and the raw
coordinate has no finite derivative through this double.
This example is in the finite polynomial heat class. It
exhibits the limitation of the split labels while the regular
integer moments remain analytic.
